\lysh{34448}{4}
\begin{alist}
\item 

$f_1 = 0,125$ άρα $f_1\% = 12,5$. \\
$f_2 = 0,275$ άρα $f_2\% = 27,5$. \\
$f_3 = 0,3$ άρα $f_3\% = 30$. \\
$f_4 = 0,075$ άρα $f_4\% = 7,5$. \\
$f_5 = 0,225$ άρα $f_5\% = 22,5$. \\
\\
$N_1 = \nu_1 = 50$. \\
$N_2 = N_1 + \nu_2 = 50 + 110 = 160$. \\
$N_3 = N_2 + \nu_3 = 160 + 120 = 280$. \\
$N_4 = N_3 + \nu_4 = 280 + 30 = 310$. \\
$N_5 = N_4 + \nu_5 = 310 + 90 = 400$ ή $N_5 = \nu = 400$. \\
\\
$F_1 = f_1 = 0,125$. \\
$F_2 = F_1 + f_2 = 0,125 + 0,275 = 0,4$. \\
$F_3 = F_2 + f_3 = 0,4 + 0,3 = 0,7$. \\
$F_4 = F_3 + f_4 = 0,7 + 0,075 = 0,775$. \\
$F_5 = F_4 + f_5 = 0,775 + 0,225 = 1$. \\
\\
$F_1 = 0,125$ άρα $F_1\% = 12,5$. \\
$F_2 = 0,4$ άρα $F_2\% = 40$. \\
$F_3 = 0,7$ άρα $F_3\% = 70$. \\
$F_4 = 0,775$ άρα $F_4\% = 77,5$. \\
$F_5 = 1$ άρα $F_5\% = 100$.

\item 
\begin{center}
\begin{mytblr}{}

{Αριθμός \\ επιβατών $x_i$} & {Αριθμός \\ αυτοκινήτων $\nu_i$} & $f_i$ & $f_i\%$ & $N_i$ & $F_i$ & $F_i\%$ \\

1 & 50 & 0,125 & 12,5 & 50 & 0,125 & 12,5 \\

2 & 110 & 0,275 & 27,5 & 160 & 0,4 & 40 \\

3 & 120 & 0,3 & 30 & 280 & 0,7 & 70 \\

4 & 30 & 0,075 & 7,5 & 310 & 0,775 & 77,5 \\

5 & 90 & 0,225 & 22,5 & 400 & 1 & 100 \\

\text{ΣΥΝΟΛΑ} & $\nu=400$ & 1 & 100 & & & \\

\end{mytblr}
\end{center}

\begin{center}
\begin{mytblr}{}

{Αριθμός \\ επιβατών $x_i$} & {Αριθμός \\ αυτοκινήτων $\nu_i$} & $x_i \nu_i$ & $(x_i - \bar{\chi})^2 \nu_i$ \\

1 & 50 & 50 & 200 \\

2 & 110 & 220 & 110 \\

3 & 120 & 360 & 0 \\

4 & 30 & 120 & 30 \\

5 & 90 & 450 & 360 \\

\text{ΣΥΝΟΛΑ} & $\nu=400$ & 1200 & 700 \\

\end{mytblr}
\end{center}

Η μέση τιμή είναι:
\[
\bar{\chi} = \frac{\sum_{i=1}^{5} x_i \nu_i}{\sum_{i=1}^{5} \nu_i} = \frac{\chi_1 \nu_1 + \chi_2 \nu_2 + \chi_3 \nu_3 + \chi_4 \nu_4 + \chi_5 \nu_5}{\nu} = \frac{50 + 220 + 360 + 120 + 450}{400} = \frac{1200}{400} = 3
\]
Αφού οι παρατηρήσεις του δείγματος είναι $\nu=400$ (άρτιος), η διάμεσος θα ίση με το ημιάθροισμα των δύο μεσαίων παρατηρήσεων, αν αυτές έχουν διαταχθεί κατά αύξουσα σειρά. Δηλαδή:
\[
\delta = \frac{\chi_{200} + \chi_{201}}{2} = \frac{3+3}{2} = 3
\]
\item  Για να ελέγξουμε την ομοιογένεια του δείγματος θα πρέπει να υπολογίσουμε το συντελεστή μεταβλητότητας $CV$ και συνεπώς την τυπική απόκλιση. Γνωρίζουμε ότι ο τύπος που μας δίνει τη διακύμανση είναι:
\[
s^2 = \frac{1}{\nu} \sum_{i=1}^{5} (x_i - \bar{\chi})^2 \nu_i \text{, άρα:}
\]
\[
s^2 = \frac{(\chi_1 - \bar{\chi})^2 \nu_1 + (\chi_2 - \bar{\chi})^2 \nu_2 + (\chi_3 - \bar{\chi})^2 \nu_3 + (\chi_4 - \bar{\chi})^2 \nu_4 + (\chi_5 - \bar{\chi})^2 \nu_5}{400} = \frac{200 + 110 + 0 + 30 + 360}{400} = \frac{700}{400} = \frac{7}{4}
\]
Οπότε η τυπική απόκλιση είναι $s = \sqrt{\frac{7}{4}} = \frac{\sqrt{7}}{2} \approx \frac{2,65}{2} = 1,32$.

Ο συντελεστής μεταβλητότητας είναι:
\[
CV = \frac{s}{\bar{\chi}} = \frac{\sqrt{7}/2}{3} = \frac{\sqrt{7}}{6} \approx \frac{2,65}{6} \approx 0,44 > \frac{1}{10}
\]
Αφού $CV > 10\%$ το δείγμα σημαίνει ότι δεν είναι ομοιογενές, συνεπώς ο ισχυρισμός του Γιάννη είναι λανθασμένος.
\end{alist}
\vspace{6mm}
\noindent\rule{\textwidth}{0.4pt}
\vspace{6mm}

